% Exercice: Exercice 4
% Sujet: Algèbre
% Généré automatiquement

\begin{exercice}[Algèbre]
\label{ex:6}

EXERCICE 4\\
$(4,2^{5} points)$\\
Le plan est muni dun repère orthonormé $(0 ,1 ,J)$. La\\
courbe représentative (Cr) ci-contre est celle dune\\
fonction f , numérique à variable réelle.\\
En $exploitant cette courbe; déterminer$\\
\item a) L'ensemble de définition de f
0,2^\{5\}pt\\
\item b) Les $limites de f en$
0,5pt\\
c et en\\
1-\\
\item c) Une équation de l'asymptote verticale à la
0,2^\{5\}pt\\
courbe\\
f(-2), f (0) et f' (0).\\
0,7^\{5\}pt\\
\item e) Le sens des variations de
~2[ 0,5pt\\
sur ] -0\\
\item 2. On suppose que la fonction f est définie pour
tout x de son ensemble de définition par\\
f(x)\\
ax + b +\\
1+1\\
\item a) En $exploitant la question 1.d) , montrer$
1pt\\
$que a = 1$,$ b = 0 et c$\\
=1\\
\item b) Montrer que la droite (D) d'é$quation y = x est asymptote oblique $à
(c;)\\
0,5pt\\
\item 3. Calculer f'(x) pour tout x de l'ensemble de définition de f
0,5pt\\

\end{exercice}

% --- Fin Exercice 4 ---
